\documentclass{ximera}

\title{u-Substitution Review Activity Facilitator Guide}
\author{MATH 426: Calculus II}

\begin{document}
\begin{abstract}
   Working with the peers in your group, solve the following problems. Make sure to show and justify all your work. Make sure everyone in the group understands the solution and participates. Be prepared to report your answers to the whole class. 
\end{abstract}
\maketitle

\textbf{Student Learning Outcomes}\\
Students should be able to:
\begin{itemize}
    \item Discuss and evaluate methods for evaluating integrals from class so far ($u$-sub and FTC).
    \item Create examples of functions based on the method needed to solve them ($u$-sub and FTC).
    \item Evaluate (indefinite and definite) integrals using $u$-sub. 
\end{itemize}

\textbf{Prerequisite Knowledge:} Students should be familiar with the Fundamental Theorem of Calculus, and $u$-substitution (perfect, imperfect and linear/affine substitution).\\
This is meant to be a more ``advanced'' $u$-sub activity, to follow an introductory activity.


The start of this activity is more open-ended and conceptual than the other activities we have seen so far in the course. As we move forward with the content, you will be exposed to several independent techniques of integration. To fully master this material you will need to be able to \textbf{choose} appropriate techniques for problems `in the wild.' This choice is \textbf{important} and it is not something that we show in our written work. In our first instantiation of this activity you will be asked to discuss your strategy for selecting an integration technique, and to share your insights with your peers. If you can begin to articulate what you `see' in mathematical problems you will be better prepared to choose the correct procedures for handling new problems.\\

\begin{exercise}
    Discuss the different integration techniques we have discussed so far in Calculus 2 (FTC and $u$-substitution). (Consider the following questions as discussion starters, you need not address all the questions in your discussion, but try to hash out any insight or confusion using these questions as a guide.)\\
    \textcolor{blue}{We will not provide solutions to these questions, as in many cases, the answers will be personal to you.  However, if you are unsure of your answers, check in with your instructor, TA or the MAC to confirm your ideas match the mathematics.}

    \begin{itemize}
        \item What makes an integration problem easy (or hard)?

        \item What features of the integrand make you think of specific techniques?

        \item If you are unsure how to proceed on a problem, what do you do next?

        \item Do you treat definite and indefinite integrals differently? How?

        \item What kinds of problems can you do in your head?
        
        \item After your discussion, compress your thoughts into a single piece of advice which is helpful for you:
    \end{itemize}
    \textcolor{orange}{Please make sure to check in with groups about their answers to some of these questions, especially the last one.  We are trying to have groups think independently and use some metacognition to think about how they are solving problems, but we also want to make sure their advice to themselves is sound! }
\end{exercise}

\begin{exercise}
    As a group create one example problem that you think is the perfect example of each of the following situations: (You can create group examples, or choose your own personal examples. You do not need to complete the problems here, but it may be helpful to annotate the problems by adding notes showing the features you see.)
        \begin{enumerate}
            \item An integral which relies on the fundamental theorem of calculus.\\
            \textcolor{blue}{ Example: $\int_0^4 3x^2dx$}

            \item An integral which has a perfect $u$-substitution.\\
            \textcolor{blue}{ Example: $\int5x^4\sin(x^5)dx$}

            \item An integral which has an {\bf{imperfect}} $u$-substitution.\\
            \textcolor{blue}{ Example: $\int x^3\sqrt{x^2+1}dx$}

            \item An integral which requires an {\bf{affine}} substitution to be tractable.\\
            \textcolor{blue}{ Example: $\int \sqrt{3x-1}dx$}
        \end{enumerate}
\end{exercise}

\begin{exercise}
    Evaluate the following integrals.
    \begin{enumerate}
        \item $\int \cot(x)dx $ (hint: you will need to rewrite $\cot(x)$ in order to be able to complete this integral.)\\
        \textcolor{blue}{$=\int \frac{\cos(x)}{\sin(x)}$\\
        We want our $\frac{du}{dx}$ to show up in the numerator, so we will choose $u=\sin(x)$.  Then $\frac{du}{dx}=\cos(x)$.\\
        $=\int \frac{\frac{du}{dx}}{u}dx=\int \frac{1}{u}du=\ln|u|+c=\ln|\sin(x)|+c$}
        
        \item $\int_0^1 \frac{e^x}{5+e^x}dx$\\
        \textcolor{blue}{Let $u=e^x$, then $\frac{du}{dx}=e^x$.  Our endpoints become $u(0)=e^0=1$ and $u(1)=e^1=e$\\
        $=\int_{u=1}^{u=e} \frac{\frac{du}{dx}}{5+u}dx=\int_1^e \frac{1}{5+u}du=\ln|5+u|\bigg|_1^e=\ln(5+e)-\ln(6)$}
        
        \item $\int x^2(x^3+4)^6dx$\\
        \textcolor{blue}{Let $u=x^3+4$, then $\frac{du}{dx}=3x^2$ or $\frac{1}{3}\frac{du}{dx}=x^2$\\
        $=\int \frac{1}{3}\frac{du}{dx}u^6dx=\int \frac{u^6}{3}du=\frac{1}{3}\frac{u^7}{7}+c=\frac{(x^3+4)^7}{21}+c$}
        
        \item $\int \frac{(\ln(x)^5)}{x}dx$\\
        \textcolor{blue}{Let $u=\ln(x)$, then $\frac{du}{dx}=\frac{1}{x}$\\
        $=\int u^5\frac{du}{dx}dx=\int u^5 du=\frac{u^6}{6}+c=\frac{(\ln(x))^6}{6}+c$}
        
        \item $\int_{-1}^2 2x^3(x^4+3)dx$\\
        \textcolor{blue}{There are multiple ways to think of this integral, but both should come out the same:\\
        $u$-sub: Let $u=x^4+3$, then $\frac{du}{dx}=4x^3$, so $\frac{1}{2}\frac{du}{dx}=2x^3$.  Our endpoints become $u(-1)=1+3=4$ and $u(2)=16+3=19$.\\
        $=\int_{u=4}^{u=19}\frac{1}{2}\frac{du}{dx} u dx=\int_4^{19} \frac12 udu=\frac{u^2}{4}\bigg|_4^{19}=\frac{19^2}{4}-\frac{4^2}{4}=\frac{345}{4}$\\
        By mulitplying through first:\\
        $=\int_{-1}^{2} 2x^7+6x^3 dx=\frac{2x^8}{8}+\frac{6x^4}{4}\bigg|_{-1}^2=\frac{2^8}{4}+\frac{3(2^4)}{2}-\frac{(-1)^8}{4}-\frac{3(-1)^4}{2}=\frac{345}{4}$}
    \end{enumerate}
\end{exercise}

\begin{exercise}
    Evalute the integrals (a) $\int x\sqrt{x-1}dx$, (b) $\int x\sqrt{x^2-1}dx$ and (c) $\int x^2\sqrt{x-1}dx$.  How are the methods for each of these integrals the same?  How are they different?  What characteristics caused you to choose your method for each integral?\\
    \textcolor{blue}{(a): Let $u=x-1$, then $\frac{du}{dx}=1$.  Note also that $x=u+1$.\\
    $=\int (u+1)\sqrt{u}\frac{du}{dx}dx=\int (u+1)u^{1/2}du=\int u^{3/2}+u^{1/2}=\frac{2u^{5/2}}{5}+\frac{2u^{3/2}}{3}+c=\frac{2(x-1)^{5/2}}{5}+\frac{2(x-1)^{3/2}}{3}+c$\\
    (b): Let $u=x^2-1$, then $\frac{du}{dx}=2x$, or $\frac{1}{2}\frac{du}{dx}=x$.\\
    $=\int \frac{1}{2}\frac{du}{dx}\sqrt{u}dx=\int \frac{1}{2}u^{1/2}du=\frac{u^{3/2}}{3}+c=\frac{(x^2-1)^{3/2}}{3}+c$\\
    (c): Let $u=x-1$, then $\frac{du}{dx}=1$.  Note that $x^2=(u+1)^2=u^2+2u+1$\\
    $=\int (u^2+2u+1)\sqrt{u}\frac{du}{dx}dx=\int (u^2+2u+1)u^{1/2}du=\int u^{5/2}+2u^{3/2}+u^{1/2}du$\\
    $=\frac{2u^{7/2}}{7}+\frac{u^{5/2}}{5}+\frac{2u^{3/2}}{3}+c=\frac{2(x-1)^{7/2}}{7}+\frac{(x-1)^{5/2}}{5}+\frac{2(x-1)^{3/2}}{3}+c$\\
    The similarity is that the $u$ is always the part inside the radical.  Because we can't have the $\frac{du}{dx}$ inside the radical, it doesn't matter when the $x^2$ is outside the radical.  (a) and (c) have the exam $u$, and we have to ``insert'' $\frac{du}{dx}$ for 1. \\
    The difference is that in (b), we have an imperfect $u$-sub, while (a) and (c) require us to both insert $\frac{du}{dx}$ and then also sub back in for another $x$ in the equation.  }\\
    \textcolor{orange}{These are not trivial integrals, and this question may take groups a while. If you are short on time, have students choose their $u$ and find $\frac{du}{dx}$, then do the substitution.  That should be enough for them to start seeing the patterns.  (c) will likely be the hardest for students, and some will try to use $u=x^2$. }
\end{exercise}



\end{document} 